\section{Locally quasi-compact spaces}
\label{section-locally-quasi-compact}

\noindent
Recall that a neighbourhood of a point need not be open.

\begin{definition}
\label{definition-locally-quasi-compact}
A topological space $X$ is called
{\it locally quasi-compact}\footnote{This may not be standard notation.
Alternative notions used in the literature are: (1) Every point has some
quasi-compact neighbourhood, and (2) Every point has a closed quasi-compact
neighbourhood. A scheme has the property that every point has a fundamental
system of open quasi-compact neighbourhoods.} if every
point has a fundamental system of quasi-compact neighbourhoods.
\end{definition}

\noindent
The term {\it locally compact space} in the literature often refers
to a space as in the following lemma.

\begin{lemma}
\label{lemma-locally-quasi-compact-Hausdorff}
A Hausdorff space is locally quasi-compact if and only if every point
has a quasi-compact neighbourhood.
\end{lemma}

\begin{proof}
Let $X$ be a Hausdorff space. Let $x \in X$ and let $x \in E \subset X$
be a quasi-compact neighbourhood. Then $E$ is closed by
Lemma \ref{lemma-quasi-compact-in-Hausdorff}.
Suppose that $x \in U \subset X$ is an open neighbourhood of $x$.
Then $Z = E \setminus U$ is a closed subset of $E$ not containing $x$.
Hence we can find a pair of disjoint open subsets $W, V \subset E$
of $E$ such that $x \in V$ and $Z \subset W$, see
Lemma \ref{lemma-quasi-compact-in-Hausdorff}.
It follows that $\overline{V} \subset E$ is a closed neighbourhood
of $x$ contained in $E \cap U$. Also $\overline{V}$ is quasi-compact
as a closed subset of $E$ (Lemma \ref{lemma-closed-in-quasi-compact}).
In this way we obtain a fundamental system of quasi-compact neighbourhoods
of $x$.
\end{proof}

\begin{lemma}[Baire category theorem]
\label{lemma-baire-category-locally-compact}
Let $X$ be a locally quasi-compact Hausdorff space.
Let $U_n \subset X$, $n \geq 1$ be dense open subsets. Then
$\bigcap_{n \geq 1} U_n$ is dense in $X$.
\end{lemma}

\begin{proof}
After replacing $U_n$ by $\bigcap_{i = 1, \ldots, n} U_i$
we may assume that $U_1 \supset U_2 \supset \ldots$.
Let $x \in X$. We will show that $x$ is in the closure of
$\bigcap_{n \geq 1} U_n$. Thus let $E$ be a neighbourhood of $x$.
To show that $E \cap \bigcap_{n \geq 1} U_n$ is nonempty we
may replace $E$ by a smaller neighbourhood. After replacing
$E$ by a smaller neighbourhood, we may assume that $E$ is quasi-compact.

\medskip\noindent
Set $x_0 = x$ and $E_0 = E$. Below, we will inductively choose
a point $x_i \in E_{i - 1} \cap U_i$ and a quasi-compact
neighbourhood $E_i$ of $x_i$ with $E_i \subset E_{i - 1} \cap U_i$.
Because $X$ is Hausdorff, the subsets $E_i \subset X$ are closed
(Lemma \ref{lemma-quasi-compact-in-Hausdorff}).
Since the $E_i$ are also nonempty we conclude that
$\bigcap_{i \geq 1} E_i$ is nonempty
(Lemma \ref{lemma-intersection-closed-in-quasi-compact}).
Since $\bigcap_{i \geq 1} E_i \subset E \cap \bigcap_{n \geq 1} U_n$
this proves the lemma.

\medskip\noindent
The base case $i = 0$ we have done above. Induction step.
Since $E_{i - 1}$ is a neighbourhood of $x_{i - 1}$ we can
find an open $x_{i - 1} \in W \subset E_{i - 1}$.
Since $U_i$ is dense in $X$
we see that $W \cap U_i$ is nonempty.
Pick any $x_i \in W \cap U_i$.
By definition of locally quasi-compact spaces we can
find a quasi-compact neighbourhood $E_i$ of $x_i$
contained in $W \cap U_i$. Then $E_i \subset E_{i - 1} \cap U_i$
as desired.
\end{proof}

\begin{lemma}
\label{lemma-relatively-compact-refinement}
Let $X$ be a Hausdorff and quasi-compact space.
Let $X = \bigcup_{i \in I} U_i$ be an open covering.
Then there exists an open covering $X = \bigcup_{i \in I} V_i$
such that $\overline{V_i} \subset U_i$ for all $i$.
\end{lemma}

\begin{proof}
Let $x \in X$. Choose an $i(x) \in I$ such that $x \in U_{i(x)}$.
Since $X \setminus U_{i(x)}$ and $\{x\}$ are disjoint closed
subsets of $X$, by Lemmas \ref{lemma-closed-in-quasi-compact} and
\ref{lemma-quasi-compact-in-Hausdorff}
there exists an open neighbourhood $U_x$ of $x$
whose closure is disjoint from $X \setminus U_{i(x)}$.
Thus $\overline{U_x} \subset U_{i(x)}$. Since $X$ is quasi-compact,
there is a finite list of points $x_1, \ldots, x_m$ such that
$X = U_{x_1} \cup \ldots \cup U_{x_m}$. Setting
$V_i = \bigcup_{i = i(x_j)} U_{x_j}$ the proof is finished.
\end{proof}

\begin{lemma}
\label{lemma-refine-covering}
Let $X$ be a Hausdorff and quasi-compact space.
Let $X = \bigcup_{i \in I} U_i$ be an open covering.
Suppose given an integer $p \geq 0$ and for every $(p + 1)$-tuple
$i_0, \ldots, i_p$ of $I$ an open covering
$U_{i_0} \cap \ldots \cap U_{i_p} = \bigcup W_{i_0 \ldots i_p, k}$.
Then there exists an open covering $X = \bigcup_{j \in J} V_j$
and a map $\alpha : J \to I$ such that $\overline{V_j} \subset U_{\alpha(j)}$
and such that each $V_{j_0} \cap \ldots \cap V_{j_p}$
is contained in $W_{\alpha(j_0) \ldots \alpha(j_p), k}$
for some $k$.
\end{lemma}

\begin{proof}
Since $X$ is quasi-compact, there is a reduction to the
case where $I$ is finite (details omitted).
We prove the result for $I$ finite by induction on $p$.
The base case $p = 0$ is immediate by taking a covering as in
Lemma \ref{lemma-relatively-compact-refinement}
refining the open covering $X = \bigcup W_{i_0, k}$.

\medskip\noindent
Induction step. Assume the lemma proven for $p - 1$.
For all $p$-tuples $i'_0, \ldots, i'_{p - 1}$ of $I$ let
$U_{i'_0} \cap \ldots \cap U_{i'_{p - 1}} =
\bigcup W_{i'_0 \ldots i'_{p - 1}, k}$
be a common refinement of the coverings
$U_{i_0} \cap \ldots \cap U_{i_p} = \bigcup W_{i_0 \ldots i_p, k}$
for those $(p + 1)$-tuples such that
$\{i'_0, \ldots, i'_{p - 1}\} = \{i_0, \ldots, i_p\}$ (equality of sets).
(There are finitely many of these as $I$ is finite.)
By induction there exists a solution for these opens, say
$X = \bigcup V_j$ and $\alpha : J \to I$.
At this point the covering $X = \bigcup_{j \in J} V_j$
and $\alpha$ satisfy $\overline{V_j} \subset U_{\alpha(j)}$
and each $V_{j_0} \cap \ldots \cap V_{j_p}$
is contained in $W_{\alpha(j_0) \ldots \alpha(j_p), k}$
for some $k$ if there is a repetition in $\alpha(j_0), \ldots, \alpha(j_p)$.
Of course, we may and do assume that $J$ is finite.

\medskip\noindent
Fix $i_0, \ldots, i_p \in I$ pairwise distinct. Consider $(p + 1)$-tuples
$j_0, \ldots, j_p \in J$ with $i_0 = \alpha(j_0), \ldots, i_p = \alpha(j_p)$
such that $V_{j_0} \cap \ldots \cap V_{j_p}$
is {\bf not} contained in $W_{\alpha(j_0) \ldots \alpha(j_p), k}$ for any $k$.
Let $N$ be the number of such $(p + 1)$-tuples. We will show how to decrease
$N$. Since
$$
\overline{V_{j_0}} \cap \ldots \cap \overline{V_{j_p}} \subset
U_{i_0} \cap \ldots \cap U_{i_p} = \bigcup W_{i_0 \ldots i_p, k}
$$
we find a finite set $K = \{k_1, \ldots, k_t\}$ such that the LHS
is contained in $\bigcup_{k \in K} W_{i_0 \ldots i_p, k}$.
Then we consider the open covering
$$
V_{j_0} =
(V_{j_0} \setminus (\overline{V_{j_1}} \cap \ldots \cap \overline{V_{j_p}}))
\cup (\bigcup\nolimits_{k \in K} V_{j_0} \cap W_{i_0 \ldots i_p, k})
$$
The first open on the RHS intersects $V_{j_1} \cap \ldots \cap V_{j_p}$
in the empty set and the other opens $V_{j_0, k}$ of the RHS
satisfy $V_{j_0, k} \cap V_{j_1} \ldots \cap V_{j_p} \subset
W_{\alpha(j_0) \ldots \alpha(j_p), k}$.
Set $J' = J \amalg K$. For $j \in J$ set $V'_j = V_j$ if $j \not = j_0$
and set $V'_{j_0} =
V_{j_0} \setminus (\overline{V_{j_1}} \cap \ldots \cap \overline{V_{j_p}})$.
For $k \in K$ set $V'_k = V_{j_0, k}$. Finally, the map $\alpha' : J' \to I$
is given by $\alpha$ on $J$ and maps every element of $K$ to $i_0$.
A simple check shows that $N$ has decreased by one under this replacement.
Repeating this procedure $N$ times we arrive at the situation where
$N = 0$.

\medskip\noindent
To finish the proof we argue by induction on the number $M$ of $(p + 1)$-tuples
$i_0, \ldots, i_p \in I$ with pairwise distinct entries for which there exists
a $(p + 1)$-tuple $j_0, \ldots, j_p \in J$ with
$i_0 = \alpha(j_0), \ldots, i_p = \alpha(j_p)$ such that
$V_{j_0} \cap \ldots \cap V_{j_p}$ is {\bf not} contained in
$W_{\alpha(j_0) \ldots \alpha(j_p), k}$ for any $k$. To do this, we
claim that the operation performed in the previous paragraph does not
increase $M$. This follows formally from the fact that the map
$\alpha' : J' \to I$ factors through a map $\beta : J' \to J$
such that $V'_{j'} \subset V_{\beta(j')}$.
\end{proof}

\begin{lemma}
\label{lemma-lift-covering-of-a-closed}
Let $X$ be a Hausdorff and locally quasi-compact space.
Let $Z \subset X$ be a quasi-compact (hence closed) subset.
Suppose given an integer $p \geq 0$, a set $I$, for every $i \in I$
an open $U_i \subset X$, and for every $(p + 1)$-tuple
$i_0, \ldots, i_p$ of $I$ an open
$W_{i_0 \ldots i_p} \subset U_{i_0} \cap \ldots \cap U_{i_p}$
such that
\begin{enumerate}
\item $Z \subset \bigcup U_i$, and
\item for every $i_0, \ldots, i_p$ we have
$W_{i_0 \ldots i_p} \cap Z = U_{i_0} \cap \ldots \cap U_{i_p} \cap Z$.
\end{enumerate}
Then there exist opens $V_i$ of $X$ such that
we have $Z \subset \bigcup V_i$,
for all $i$ we have $\overline{V_i} \subset U_i$, and
we have $V_{i_0} \cap \ldots \cap V_{i_p} \subset W_{i_0 \ldots i_p}$
for all $(p + 1)$-tuples $i_0, \ldots, i_p$.
\end{lemma}

\begin{proof}
Since $Z$ is quasi-compact, there is a reduction to the
case where $I$ is finite (details omitted).
Because $X$ is locally quasi-compact and $Z$ is quasi-compact,
we can find a neighbourhood $Z \subset E$ which is quasi-compact,
i.e., $E$ is quasi-compact and contains an open neighbourhood
of $Z$ in $X$. If we prove the result after replacing $X$ by $E$,
then the result follows. Hence we may assume $X$ is quasi-compact.

\medskip\noindent
We prove the result in case $I$ is finite and $X$ is quasi-compact
by induction on $p$. The base case is $p = 0$. In this case we have
$X = (X \setminus Z) \cup \bigcup W_i$. By
Lemma \ref{lemma-relatively-compact-refinement}
we can find a covering $X = V \cup \bigcup V_i$ by
opens $V_i \subset W_i$ and $V \subset X \setminus Z$
with $\overline{V_i} \subset W_i$ for all $i$. Then we see that
we obtain a solution of the problem posed by the lemma.

\medskip\noindent
Induction step. Assume the lemma proven for $p - 1$.
Set $W_{j_0 \ldots j_{p - 1}}$ equal to the intersection of
all $W_{i_0 \ldots i_p}$ with
$\{j_0, \ldots, j_{p - 1}\} = \{i_0, \ldots, i_p\}$ (equality of sets).
By induction there exists a solution for these opens, say
$V_i \subset U_i$.
It follows from our choice of $W_{j_0 \ldots j_{p - 1}}$ that we have
$V_{i_0} \cap \ldots \cap V_{i_p} \subset W_{i_0 \ldots i_p}$
for all $(p + 1)$-tuples $i_0, \ldots, i_p$ where $i_a = i_b$ for
some $0 \leq a < b \leq p$.
Thus we only need to modify our choice of
$V_i$ if $V_{i_0} \cap \ldots \cap V_{i_p} \not \subset W_{i_0 \ldots i_p}$
for some $(p + 1)$-tuple $i_0, \ldots, i_p$ with pairwise distinct elements.
In this case we have
$$
T =
\overline{V_{i_0} \cap \ldots \cap V_{i_p} \setminus W_{i_0 \ldots i_p}}
\subset 
\overline{V_{i_0}} \cap \ldots \cap \overline{V_{i_p}} \setminus
W_{i_0 \ldots i_p}
$$
is a closed subset of $X$ contained in $U_{i_0} \cap \ldots \cap U_{i_p}$
not meeting $Z$. Hence we can replace $V_{i_0}$ by $V_{i_0} \setminus T$
to ``fix'' the problem. After repeating this finitely many times for each
of the problem tuples, the lemma is proven.
\end{proof}

\begin{lemma}
\label{lemma-lift-covering-of-quasi-compact-hausdorff-subset}
Let $X$ be a topological space. Let $Z \subset X$ be a quasi-compact subset
such that any two points of $Z$ have disjoint open neighbourhoods in $X$.
Suppose given an integer $p \geq 0$, a set $I$, for every $i \in I$
an open $U_i \subset X$, and for every $(p + 1)$-tuple
$i_0, \ldots, i_p$ of $I$ an open
$W_{i_0 \ldots i_p} \subset U_{i_0} \cap \ldots \cap U_{i_p}$
such that
\begin{enumerate}
\item $Z \subset \bigcup U_i$, and
\item for every $i_0, \ldots, i_p$ we have
$W_{i_0 \ldots i_p} \cap Z = U_{i_0} \cap \ldots \cap U_{i_p} \cap Z$.
\end{enumerate}
Then there exist opens $V_i$ of $X$ such that
\begin{enumerate}
\item $Z \subset \bigcup V_i$,
\item $V_i \subset U_i$ for all $i$,
\item $\overline{V_i} \cap Z \subset U_i$ for all $i$, and
\item $V_{i_0} \cap \ldots \cap V_{i_p} \subset W_{i_0 \ldots i_p}$
for all $(p + 1)$-tuples $i_0, \ldots, i_p$.
\end{enumerate}
\end{lemma}

\begin{proof}
Since $Z$ is quasi-compact, there is a reduction to the
case where $I$ is finite (details omitted).
We prove the result in case $I$ is finite by induction on $p$.

\medskip\noindent
The base case is $p = 0$.
For $z \in Z \cap U_i$ and $z' \in Z \setminus U_i$ there exist
disjoint opens $z \in V_{z, z'}$ and $z' \in W_{z, z'}$ of $X$.
Since $Z \setminus U_i$ is quasi-compact
(Lemma \ref{lemma-closed-in-quasi-compact}),
we can choose a finite number $z'_1, \ldots, z'_r$ such that
$Z \setminus U_i \subset W_{z, z'_1} \cup \ldots \cup W_{z, z'_r}$.
Then we see that
$V_z = V_{z, z'_1} \cap \ldots \cap V_{z, z'_r} \cap U_i$
is an open neighbourhood of $z$ contained in $U_i$
with the property that $\overline{V_z} \cap Z \subset U_i$.
Since $z$ and $i$ where arbitrary and since $Z$ is quasi-compact
we can find a finite list $z_1, i_1, \ldots, z_t, i_t$
and opens $V_{z_j} \subset U_{i_j}$ with
$\overline{V_{z_j}} \cap Z \subset U_{i_j}$
and $Z \subset \bigcup V_{z_j}$.
Then we can set $V_i = W_i \cap (\bigcup_{j : i = i_j} V_{z_j})$
to solve the problem for $p = 0$.

\medskip\noindent
Induction step. Assume the lemma proven for $p - 1$.
Set $W_{j_0 \ldots j_{p - 1}}$ equal to the intersection of
all $W_{i_0 \ldots i_p}$ with
$\{j_0, \ldots, j_{p - 1}\} = \{i_0, \ldots, i_p\}$ (equality of sets).
By induction there exists a solution for these opens, say
$V_i \subset U_i$.
It follows from our choice of $W_{j_0 \ldots j_{p - 1}}$ that we have
$V_{i_0} \cap \ldots \cap V_{i_p} \subset W_{i_0 \ldots i_p}$
for all $(p + 1)$-tuples $i_0, \ldots, i_p$ where $i_a = i_b$ for
some $0 \leq a < b \leq p$.
Thus we only need to modify our choice of
$V_i$ if $V_{i_0} \cap \ldots \cap V_{i_p} \not \subset W_{i_0 \ldots i_p}$
for some $(p + 1)$-tuple $i_0, \ldots, i_p$ with pairwise distinct elements.
In this case we have
$$
T =
\overline{V_{i_0} \cap \ldots \cap V_{i_p} \setminus W_{i_0 \ldots i_p}}
\subset
\overline{V_{i_0}} \cap \ldots \cap \overline{V_{i_p}} \setminus
W_{i_0 \ldots i_p}
$$
is a closed subset of $X$ not meeting $Z$ by our property (3) of the
opens $V_i$. Hence we can replace $V_{i_0}$ by $V_{i_0} \setminus T$
to ``fix'' the problem. After repeating this finitely many times for each
of the problem tuples, the lemma is proven.
\end{proof}






